% ============================================================================================
\section{S1: Colour Mapping (Owner: Sidhanth Prabhu, BT2024027)}
\label{sec:s1}
\textbf{Method.} Every map is drawn on a Lambert conformal conic projection (standard parallels 33 and 45\,$^\circ$N), which
preserves shape over CONUS. Natural Earth 1:50m state borders and coastlines are added with Cartopy~\cite{cartopy}. Each
grid cell is rendered as an image pixel (nearest-neighbour, no smoothing), so the colour is the cell value and not an
interpolation. Cells outside the land mask show a neutral blue-grey labelled ``No data'', which no data colour uses. The
primary scalar is $\Tmax$; precipitation is the second scalar because its skewed distribution
(Table~\ref{tab:global}: 71.3\,\% zeros, median 0, maximum 401\,mm) needs a different mapping. Script:
\fpath{scivis/s1_colormapping/run.py}.

\subsection{Continuous colour maps and perceptual uniformity}
Fig.~\ref{fig:02} shows the same field (D2, 24~July, linear 5.1--51.0\,\degC) with four palettes. Fig.~\ref{fig:03} and
Table~\ref{tab:lightness} measure the CIELAB lightness $L^*$ of each palette~\cite{borland2007,crameri2020}.

\figfull{Fig_02_tmmx_continuous_colormaps.png}{Continuous colour maps of $\Tmax$ on D2 (24~Jul~2024, western heat-wave
peak), identical linear normalisation (5.1--51.0\,\degC). (a) jet, a rainbow baseline; (b) viridis; (c) inferno; (d)
cividis, which is optimised for colour-vision deficiency~\cite{nunez2018}.}{fig:02}

\figfull{Fig_03_colormap_lightness_profiles.png}{(a) CIELAB lightness $L^*$ along each colour map used in S1. (b) Each
map as a colour strip (top) and as its lightness only (bottom), i.e.\ how it reads in greyscale.}{fig:03}

\input{generated/tab_s1_lightness}

\textbf{Inferences.}
(1) The palette alone changes what appears to be structure. In jet (Fig.~\ref{fig:02}a) a sharp cyan$\rightarrow$yellow
``front'' crosses the Midwest at about 28--30\,\degC, and a second sharp edge appears at about 38\,\degC{} in the Plains.
Neither is a physical discontinuity; the temperature changes smoothly there. They are hue boundaries at jet's bright
band, where $L^*$ jumps from about 60 to 95 within 10\,\% of the range (Fig.~\ref{fig:03}a).
(2) jet hides the hottest cells. Its lightness \emph{falls} to 25 at the top end, so the 45--51\,\degC{} Central Valley,
lower Colorado River and Death Valley cells are dark red and look \emph{less} salient than the 38--42\,\degC{}
Montana--Dakota hot spot in yellow-orange. In greyscale (Fig.~\ref{fig:03}b), jet's coldest and hottest values both appear
dark.
(3) viridis, inferno and cividis have monotone $L^*$ (0 turns). inferno has the largest range ($L^*$ 0.1--97.9), so it
separates the most steps. In Fig.~\ref{fig:02}c the cool Rockies ridges (10--20\,\degC), the cool Pacific coast
(15--20\,\degC), the warm Great Lakes and New England (22--27\,\degC) and the northern-Plains heat dome (38--42\,\degC) are
all readable at once. We therefore use inferno as the default $\Tmax$ map in the rest of S1.
(4) viridis and cividis compress the 25--35\,\degC{} bulk of the East into a narrow teal or grey band. That is honest but
gives low contrast exactly where most land cells lie (the median is 31.7\,\degC{} on D2). This is the inherent cost of a
uniform ramp over a long-tailed range.
(5) turbo improves on jet's darkening at the low end but still turns in lightness (1 turn), so the same caveats apply.
Rainbow maps are only acceptable for categorical reading, not for judging magnitude.

\subsection{Discrete (quantised) colour maps}
Fig.~\ref{fig:04} compares the continuous inferno map with two \texttt{BoundaryNorm} classifications. One uses
\emph{equal-interval} 5\,\degC{} classes; the other uses \emph{quantile} classes, the deciles of the day's land cells. The
histograms below each map show where the breaks fall.

\figfull{Fig_04_tmmx_discrete_equal_vs_quantile.png}{Continuous vs.\ discrete colour mapping of $\Tmax$ on D2. (a)
Continuous; (b) ten equal 5\,\degC{} classes; (c) ten quantile (decile) classes; (d--f) land-cell histograms with the class
breaks and the share of land per class.}{fig:04}

\textbf{Inferences.}
(1) Equal-interval classes are interpretable: each boundary is an isotherm (e.g.\ the 35\,\degC{} line), so the map can be
read like a contour map. But land is spread very unevenly across classes, from 0.0\,\% to 42.7\,\% per class
(Fig.~\ref{fig:04}e). The 30--35\,\degC{} class alone covers most of the East and Southeast, which become one flat colour.
(2) Quantile classes equalise area by construction (8.9--10.8\,\% per class) and maximise contrast. However, the breaks
(24.6, 26.9, 28.6, 30.2, 31.6, 32.7, 33.7, 34.8, 37.3\,\degC) are only 1--3\,\degC{} apart in the middle and about
20\,\degC{} apart at the ends. The map therefore exaggerates differences of 1\,\degC{} across the Plains and Southeast, while
treating a 37\,\degC{} Kansas cell and a 51\,\degC{} Death Valley cell as the same class.
(3) The histogram (Fig.~\ref{fig:04}d) is skewed. A main mode at 31--34\,\degC{} (the humid East and South) has a cool
shoulder at about 25\,\degC{} (the Pacific Northwest, Great Lakes and mountains) and a hot tail of 40--51\,\degC{} (the
desert Southwest and the northern-Plains heat dome). Any classification has to choose which of these to resolve.
(4) Our recommendation: equal intervals when values must be compared with physical thresholds or across dates; quantiles
when the task is ranking regions within one map; continuous when exact gradients matter and the reader can use a colour
bar.

\subsection{Logarithmic mapping of precipitation}
Fig.~\ref{fig:05} maps precipitation on the two heaviest landfall days with a linear map scaled to the daily maximum, a
linear map clipped at 25\,mm, and a \texttt{LogNorm} over 0.1--401\,mm. Zeros cannot be log-scaled. Rather than adding an
arbitrary $\epsilon$, we draw dry land ($p=0$) as a separate cream ``dry'' category beneath the colour layer, so ``no
rain'' and ``0.1\,mm drizzle'' are never confused.

\figfull{Fig_05_pr_linear_vs_log.png}{Precipitation with linear vs.\ logarithmic colour mapping on D3 (Debby, top) and D7
(Helene, bottom). Left: linear 0--(day maximum); middle: linear clipped at 25\,mm (the extend arrow marks saturation);
right: LogNorm 0.1--401\,mm with zeros shown in cream.}{fig:05}

\textbf{Inferences.}
(1) On D3 a linear map scaled to the 331\,mm maximum shows only Debby's swath along the Florida Gulf coast and the
Georgia--South Carolina coastal plain, where totals exceed 150\,mm. Everything else looks dry, even though 32.3\,\% of land
is wet: the median wet cell receives 3.5\,mm, about 1\,\% of the colour range.
(2) Clipping at 25\,mm reveals a second, frontal rain band across the Dakotas, Minnesota, Wisconsin and Michigan and into
New England. But the hurricane core saturates, so 30\,mm and 300\,mm look identical. Clipping trades one blind spot for
another.
(3) The log map shows both at once: Debby's core (more than 100\,mm), the frontal band (5--30\,mm), and monsoon
showers of 0.3--5\,mm over Colorado, New Mexico, Utah and Arizona that are invisible in both linear versions.
(4) On D7 the log map separates Helene's core (more than 100\,mm, a SW--NE stripe over Georgia, the Carolinas and the
southern Appalachians) from a broad 10--30\,mm shield over Alabama, Tennessee and Kentucky. It also shows unrelated rain in
Washington state and New England. Only 23.4\,\% of land is wet on D7, but the median wet cell receives 12.6\,mm, much more
than on D3, because the rain is concentrated in one organised system.
(5) The limitation: equal colour steps now mean equal \emph{ratios}, so a reader must not judge differences linearly.
Going from 1 to 10\,mm is as large a colour step as going from 30 to 300\,mm.

\subsection{Sequential vs.\ diverging maps on the anomaly}
Fig.~\ref{fig:06} maps the $\Tmax$ anomaly of Eq.~\eqref{eq:anom} on D2 and D7 with two diverging maps (RdBu\_r and coolwarm
with \texttt{TwoSlopeNorm} centred at 0, limits $\pm15$\,\degC) and with sequential viridis stretched min--max.

\figfull{Fig_06_tmmx_anomaly_diverging.png}{Daily $\Tmax$ anomaly relative to each cell's Jul--Sep 2024 mean on D2 (top)
and D7 (bottom). (a, d) RdBu\_r and (b, e) coolwarm, both centred at 0 with a two-slope norm; (c, f) sequential viridis
with min--max scaling, where the black tick on the colour bar marks 0.}{fig:06}

\textbf{Inferences.}
(1) D2 is anomalously warm over 69.3\,\% of land (mean $+2.05$\,\degC). Montana, northern Wyoming and the western
Dakotas reach $+10$ to $+15$\,\degC, which is the core of the late-July heat dome. The Great Basin and Pacific Northwest
run $+3$ to $+8$\,\degC.
(2) On the same day a cool band of $-3$ to $-8$\,\degC{} runs from south Texas through Louisiana and Mississippi into
Tennessee. This is the rain-cooled Gulf moisture plume that also appears in Fig.~\ref{fig:09}b. The diverging map shows that
these two regions have opposite \emph{signs}, which a sequential map cannot convey.
(3) D7 reverses the picture: only 22.1\,\% of land is warmer than usual (mean $-2.48$\,\degC), and 20.0\,\% is more than
5\,\degC{} colder. The strongest cold anomaly ($-10$ to $-13$\,\degC) is a narrow north--south band over Alabama, Georgia
and Tennessee directly under Helene's cloud and rain shield, while the northern Plains are $+5$ to $+8$\,\degC. So one
event depresses daytime maximum temperature locally by as much as the whole seasonal cooling from July to September.
(4) viridis has no visual zero. On D2 the value 0\,\degC{} falls on a dark teal, so a slightly warm and a slightly cool cell
look alike, and ``yellow'' only says ``highest anomaly of this map''. Diverging maps are right when the data have a
meaningful midpoint (an anomaly, a difference, a correlation). Sequential maps are right when the data are a magnitude
with a natural zero, such as temperature, VPD or rain~\cite{moreland2009}.
(5) RdBu\_r saturates more strongly than coolwarm. coolwarm's lighter extremes ($L^*\geq38$, Table~\ref{tab:lightness})
keep state borders legible, at the price of weaker emphasis on the $\pm10$\,\degC{} extremes.

\subsection{Global vs.\ local normalisation across timesteps}
Figs.~\ref{fig:07} and \ref{fig:08} show $\Tmax$ on all eight dates. Fig.~\ref{fig:07} uses one \emph{global} colour range,
the 92-day minimum and maximum ($-6.35$ to 53.85\,\degC). Fig.~\ref{fig:08} stretches each date to its own
\emph{local} minimum and maximum. Fig.~\ref{fig:09} applies a global logarithmic range to precipitation.\footnote{Animated
versions (8 frames, 1 frame/s, fixed colour-bar axis, date in the title):
\fpath{images/Fig_07_tmmx_global_norm.gif}, \fpath{images/Fig_08_tmmx_local_norm.gif},
\fpath{images/Fig_09_pr_global_log.gif}.} Table~\ref{tab:s1ranges} lists the local ranges.

\figfull{Fig_07_tmmx_global_norm.png}{$\Tmax$ on the eight dates with \emph{global} normalisation: a single colour bar
spanning the Jul--Sep 2024 minimum and maximum. All eight frames of the GIF \fpath{Fig_07_tmmx_global_norm.gif}.}{fig:07}

\figfull{Fig_08_tmmx_local_norm.png}{$\Tmax$ on the eight dates with \emph{local} normalisation: each panel has its own
colour bar from that day's minimum to maximum. All eight frames of the GIF \fpath{Fig_08_tmmx_local_norm.gif}.}{fig:08}

\figfull{Fig_09_pr_global_log.png}{Precipitation on the eight dates with one global logarithmic range (0.1--401\,mm);
dry land in cream. All eight frames of the GIF \fpath{Fig_09_pr_global_log.gif}.}{fig:09}

\input{generated/tab_s1_ranges}

\textbf{Inferences.}
(1) With global normalisation the same colour means the same temperature on every date, so the seasonal trend is
visible. D1 and D2 are the brightest frames (desert Southwest yellow, Plains orange). By D6--D8 the Pacific Northwest, the
northern Rockies and New England are clearly darker. The CONUS-mean $\Tmax$ falls from 31.2\,\degC{} (D1, D2) to
26.7--27.1\,\degC{} (D6--D8) (Table~\ref{tab:perdate}).
(2) Event signatures are comparable across dates only under global scaling. The darker band over the Upper Midwest on D3
(Debby day, frontal rain) and the dark patch over Alabama, Georgia and Tennessee on D7 and D8 (Helene) are recognisable
as \emph{the same} 20--25\,\degC{} level.
(3) The cost of global scaling is within-day contrast. The 60\,\degC{} global span means most land on any day falls into
25--40\,\degC, which is only a quarter of the colour range.
(4) Local normalisation uses the full palette every day (each local range is 64--80\,\% of the global range,
Table~\ref{tab:s1ranges}), so within-day patterns are sharper. Examples are the Central Valley vs.\ the coast on D1, and
the D5 contrast between the Southwest (above 45\,\degC) and the Pacific Northwest.
(5) But local normalisation \emph{misleads} across dates. The coolest region of each day turns near-black whatever its
actual value. D6's Rockies at $-2.7$\,\degC{} and D3's Upper Midwest at about 20\,\degC{} receive the same colour. The
seasonal cooling disappears completely in the local version of the GIF, and any comparison requires reading eight
colour bars.
(6) For precipitation (Fig.~\ref{fig:09}) a global log scale is essential: a single range spans D4's quiet 100\,mm maximum
and D7's 359\,mm. The panels tell the season's story directly. D1 shows Beryl's swath from the upper Texas coast through
Arkansas and Missouri to Michigan. D2 shows Southwest monsoon storms. D3 shows Debby. D6 shows Francine's core over
southern Louisiana and Mississippi. D7 and D8 show Helene's rain moving from Florida and Georgia onto the southern
Appalachians.

\subsection{Observations and limitations (S1)}
\begin{table}[!t]
\centering
\caption{Colour-map comparison for $\Tmax$ (qualitative judgement, supported by Figs.~\ref{fig:02}--\ref{fig:06}).}
\label{tab:cmapcomp}
\footnotesize
\begin{tabular}{p{1.4cm}p{1.2cm}p{2.2cm}p{2.3cm}}
\toprule
Map & Type & Strength & Weakness \\
\midrule
jet & rainbow & many hues & false fronts, dark top end, not CVD-safe \\
viridis & sequential & uniform, CVD-safe & low contrast at mid range \\
inferno & sequential & largest $L^*$ range & dark low end can hide cool mountains \\
cividis & sequential & CVD-optimised & muted, greyish middle \\
RdBu\_r & diverging & strong sign contrast & saturated extremes \\
coolwarm & diverging & softer, borders legible & weaker extremes \\
\bottomrule
\end{tabular}
\end{table}
Table~\ref{tab:cmapcomp} summarises the palettes. The main lessons are: use perceptually uniform maps for magnitude;
use diverging maps only around a meaningful centre; use a log scale for multiplicative, zero-inflated quantities such as
rain; and use global normalisation whenever the reader compares timesteps. Local normalisation is acceptable only with a
colour bar per panel and only for within-day questions. A limitation is that gridMET $\Tmax$ is a daily maximum on a 4\,km
grid, so urban heat islands and sub-grid cold-air pools are not resolved. Also, the anomaly baseline is a single season,
not a multi-decade climatology, so ``anomalous'' here means ``unusual for summer 2024''.
